\documentclass[11pt,a4paper]{article}

% Packages for XeLaTeX
\usepackage{fontspec}
\setmainfont{Times New Roman}
\setsansfont{Arial}
\setmonofont{Consolas}

% Mathematical packages
\usepackage{amsmath}
\usepackage{amsfonts}
\usepackage{amssymb}

% Graphics and colors
\usepackage{graphicx}
\usepackage{tcolorbox}

% Layout
\usepackage[margin=2.5cm]{geometry}
\usepackage{titlesec}
\usepackage{enumitem}
\usepackage{xcolor}
\usepackage{booktabs}
\usepackage{longtable}
\usepackage{hyperref}
\usepackage{fancyhdr}
\usepackage{lastpage}

% Color definitions
\definecolor{titleblue}{RGB}{37,99,235}
\definecolor{sectiongray}{RGB}{75,85,99}

% Hyperref settings
\hypersetup{
    colorlinks=true,
    linkcolor=blue,
    filecolor=magenta,
    urlcolor=cyan,
    pdftitle={AI and Machine Learning Syllabus},
    pdfauthor={Department of Physics, Pusan National University},
}

% Title formatting
\titleformat{\section}
  {\Large\bfseries\color{titleblue}}
  {\thesection}{1em}{}[\titlerule]

\titleformat{\subsection}
  {\large\bfseries\color{sectiongray}}
  {\thesubsection}{1em}{}

% Header and footer
\pagestyle{fancy}
\fancyhf{}
\fancyhead[L]{\small AI and Machine Learning: Intersection of AI and Physics}
\fancyhead[R]{\small Spring 2026}
\fancyfoot[C]{\thepage\ / \pageref{LastPage}}

% Document starts
\begin{document}

% Title page
\begin{center}
    {\Huge\bfseries\color{titleblue} AI and Machine Learning}\\[0.2cm]
    {\Huge\bfseries\color{titleblue} Intersection of AI and Physics}\\[0.5cm]
    {\Large\textbf{From Neural Networks to Physics-Informed AI}}\\[1cm]
    
    {\large Department of Physics}\\
    {\large Pusan National University}\\[0.5cm]
    {\large Spring Semester 2026}\\[0.3cm]
    {\large Level: 2nd Year Undergraduate}\\[0.3cm]
    {\large Contact Hours: 3 hours per week (2 hours lecture + 1 hour lab)}\\[0.3cm]
    {\large Tools: Python, Claude AI, VS Code}\\[1cm]
\end{center}

% Team project info box
\noindent\fbox{\parbox{\textwidth}{
\textbf{Team Project Information}\\[0.2cm]
\begin{itemize}[leftmargin=*,noitemsep]
    \item \textbf{Team Structure}: 4 students per team
    \item \textbf{Project Format}: Collaborative team project
    \item \textbf{Presentation}: End-of-semester team presentations
    \item \textbf{Collaboration Tools}: GitHub for code sharing and version control
\end{itemize}
}}

\vspace{0.5cm}

\section{Course Description}

This course develops essential computational competencies in physics by integrating artificial intelligence and machine learning technologies in response to the demands of the AI era. The curriculum covers foundational Neural Networks, state-of-the-art Large Language Models (LLMs), and physics-specialized AI through Physics-Informed Neural Networks (PINNs).

\subsection*{Key Features}
\begin{itemize}[noitemsep]
    \item Systematic theoretical learning based on MIT course materials
    \item "Vibe coding" practice utilizing Claude AI
    \item AI-based numerical solutions for mechanics, electromagnetism, quantum mechanics, and statistical physics
    \item PINN techniques directly incorporating physical laws into learning
\end{itemize}

\section{Course Objectives}

The course pursues the following goals:

\begin{enumerate}[noitemsep]
    \item Introduction to programming methodologies in the AI era
    \item Understanding fundamental principles and operations of Neural Networks
    \item Learning Deep Learning and Transformer architectures
    \item Comprehensive understanding of AI (LLM-based) and machine learning principles
    \item Developing efficient coding skills using LLMs
    \item Application of these methods to solve problems in classical mechanics, electromagnetism, statistical physics, quantum mechanics, and related fields
    \item Gaining practical experience in solving physics problems with AI
    \item Mastering differential equation solutions using PINNs
\end{enumerate}

\section{Textbooks and References}

\subsection{Primary Textbook}
\begin{itemize}[noitemsep]
    \item \textbf{MIT 6.S191}: Introduction to Deep Learning (lecture notes and videos)
    \item Instructor-provided materials: Jupyter Notebooks, Python code examples
\end{itemize}

\subsection{Reference Books}
\begin{itemize}[noitemsep]
    \item \textit{Deep Learning} by Goodfellow, Bengio, and Courville
    \item \textit{Physics-Informed Neural Networks} (collection of related papers)
    \item \textit{Computational Physics} by Mark Newman
\end{itemize}

\subsection{Online Resources}
\begin{itemize}[noitemsep]
    \item Claude AI Documentation
    \item PyTorch/TensorFlow Tutorials
    \item ArXiv papers on PINN applications
\end{itemize}

\newpage

\section{Course Schedule (16 Weeks)}

\subsection{Part I: Neural Networks \& Deep Learning (Weeks 1-7)}

\subsubsection*{Week 1: Course Introduction and Environment Setup}
\begin{itemize}[noitemsep]
    \item Course objectives and evaluation methods
    \item Python environment setup: uv, Git, VS Code/Cursor
    \item Claude AI account creation and basic usage
    \item \textbf{Lab}: "Hello, Neural Network!" - first neural network implementation, ML basics through polynomial fitting
    \item \textbf{Key Concepts}: Development Environment, AI in Physics
\end{itemize}

\subsubsection*{Week 2: Machine Learning Fundamentals}
\begin{itemize}[noitemsep]
    \item Three types of machine learning: supervised, unsupervised, and reinforcement learning
    \item Data preprocessing and feature extraction
    \item Loss functions and optimization
    \item \textbf{Lab}: Linear Regression, clustering, Gradient Descent visualization
    \item \textbf{Key Concepts}: Supervised Learning, Loss Functions, Gradient Descent
    \item \textbf{Assignment 1}: Linear and nonlinear regression analysis using experimental data
\end{itemize}

\subsubsection*{Week 3: Neural Network Fundamentals}
\begin{itemize}[noitemsep]
    \item Perceptron and Multi-Layer Perceptron
    \item Activation Functions: ReLU, Sigmoid, Tanh
    \item Mathematical structure of Forward Propagation
    \item Universal Approximation Theorem
    \item \textbf{Lab}: Pure Neural Network implementation with Numpy
    \item \textbf{Key Concepts}: Neurons, Activation, Forward Pass
    \item \textbf{Reference}: MIT 6.S191 Lecture 1
\end{itemize}

\subsubsection*{Week 4: Learning with Physics Data}
\begin{itemize}[noitemsep]
    \item Regression models using TensorFlow/Keras
    \item Overfitting and Underfitting
    \item Relationship between model complexity and performance
    \item \textbf{Lab}: 1D function approximation, projectile motion prediction, pendulum period learning
    \item \textbf{Key Concepts}: Regression with Neural Networks, Model Complexity
    \item \textbf{Assignment 2}: Applying NN to various physics problems
\end{itemize}

\subsubsection*{Week 5: Core Deep Learning Techniques}
\begin{itemize}[noitemsep]
    \item Regularization: L1/L2, Dropout, Batch Normalization
    \item Overfitting vs. Underfitting
    \item Data Augmentation and Transfer Learning concepts
    \item \textbf{Lab}: MNIST handwriting recognition (introducing CNN)
    \item \textbf{Key Concepts}: Regularization, Generalization, CNN Basics
    \item \textbf{Reference}: MIT 6.S191 Lecture 2
\end{itemize}

\subsubsection*{Week 6: Transformers and Attention Mechanism}
\begin{itemize}[noitemsep]
    \item Limitations of RNN and emergence of Attention
    \item Principles of Self-Attention, Transformer architecture analysis
    \item Positional Encoding
    \item \textbf{Lab}: Attention mechanism, Multi-Head Attention, Transformer Block implementation
    \item \textbf{Key Concepts}: Attention, Self-Attention, Transformers, Sequence Modeling
    \item \textbf{Reference}: "Attention Is All You Need" (Vaswani et al., 2017)
\end{itemize}

\subsubsection*{Week 7: Introduction to Large Language Models (LLM)}
\begin{itemize}[noitemsep]
    \item Comparing GPT, BERT, and Claude architectures
    \item Pre-training and Fine-tuning
    \item Token, Embedding, Context Window
    \item RLHF (Reinforcement Learning from Human Feedback)
    \item \textbf{Lab}: Tokenization methods, GPT vs BERT comparison, Claude API learning
    \item \textbf{Key Concepts}: LLM Architecture, Tokenization, Prompting, RLHF
    \item \textbf{Midterm Project Presentation}: Part I learning summary and mini-project
\end{itemize}

\subsection{Part II: LLM Vibe Coding for Physics (Weeks 8-12)}

\subsubsection*{Week 8: Midterm Exam / Introduction to LLM-based Coding}
\begin{itemize}[noitemsep]
    \item What is "Vibe Coding"? - Code generation with natural language
    \item Prompt Engineering techniques
    \item Efficient debugging with Claude
    \item \textbf{Lab}: Generating physics simulation code with natural language
    \item \textbf{Key Concepts}: Prompt Engineering, AI-Assisted Programming
\end{itemize}

\subsubsection*{Week 9: Classical Mechanics Problem Solving}
\begin{itemize}[noitemsep]
    \item Numerical solutions for Newton's equations (Euler, RK4)
    \item Planetary motion simulation
    \item Pendulum motion and Chaotic Systems
    \item Lagrangian and Hamiltonian mechanics
    \item \textbf{Lab}: Euler vs RK4 comparison, planetary motion, chaotic pendulum, Lagrangian/Hamiltonian
    \item \textbf{Key Concepts}: ODEs, Numerical Integration, Chaotic Dynamics
    \item \textbf{Assignment 3}: Mechanics problem solving and visualization using LLM
\end{itemize}

\subsubsection*{Week 10: Electromagnetism Simulation}
\begin{itemize}[noitemsep]
    \item Numerical solutions for Maxwell equations (FDTD)
    \item Electric/magnetic field calculation and visualization
    \item Laplace equation and electrostatic problems
    \item Electromagnetic wave propagation simulation
    \item \textbf{Lab}: Electric/magnetic field basics, Maxwell equations 1D/2D, electromagnetic wave animation
    \item \textbf{Key Concepts}: Maxwell Equations, FDTD, Vector Field Visualization
    \item \textbf{Assignment 4}: Electric field analysis for complex charge distributions
\end{itemize}

\subsubsection*{Week 11: Quantum Mechanics Simulation}
\begin{itemize}[noitemsep]
    \item Numerical solutions for Schrödinger equation
    \item Wave function visualization
    \item Tunneling effect simulation
    \item Finite Well, Harmonic Oscillator
    \item \textbf{Lab}: Schrödinger equation, wave functions, quantum tunneling, quantum wells
    \item \textbf{Key Concepts}: Quantum Mechanics, Wave Functions, Tunneling
    \item \textbf{Assignment 5}: Quantum state analysis in various potentials
\end{itemize}

\subsubsection*{Week 12: Statistical Physics and Monte Carlo Simulation}
\begin{itemize}[noitemsep]
    \item Monte Carlo methodology
    \item Ising model and phase transitions
    \item Metropolis-Hastings algorithm
    \item Thermodynamic property calculation
    \item \textbf{Lab}: Random Walk, π estimation with Monte Carlo, 1D/2D Ising models
    \item \textbf{Key Concepts}: Statistical Physics, Monte Carlo, Phase Transitions
    \item \textbf{Midterm Project Check}: Part II lab results sharing
\end{itemize}

\subsection{Part III: Physics-Informed Neural Networks (Weeks 13-14)}

\subsubsection*{Week 13: PINN Fundamentals (ODE Edition)}
\begin{itemize}[noitemsep]
    \item Physics-Informed Neural Networks concepts
    \item Incorporating differential equations into Loss Function
    \item Automatic Differentiation
    \item PINN vs. traditional numerical methods (RK4)
    \item \textbf{Lab}: Simple ODE, harmonic oscillator, damped oscillator, Lorenz system
    \item \textbf{Key Concepts}: PINN, Physics Loss, Boundary Conditions
    \item \textbf{Reference Papers}: Raissi et al. (2019), Karniadakis et al. (2021)
\end{itemize}

\subsubsection*{Week 14: PINN Applications - Partial Differential Equations}
\begin{itemize}[noitemsep]
    \item 1D/2D Heat Equation
    \item 1D/2D Wave Equation
    \item Burgers equation (nonlinear PDE)
    \item Handling complex boundary conditions
    \item \textbf{Lab}: PINN basic structure, Heat/Wave Equation, Burgers Equation
    \item \textbf{Key Concepts}: PDEs, PINN for Spatial-Temporal Problems
    \item \textbf{Assignment 6}: PDE solution project using PINN
\end{itemize}

\subsection{Part IV: Final Project (Weeks 15-16)}

\subsubsection*{Week 15: PINN Applications II - Advanced Topics \& Final Project Start}
\begin{itemize}[noitemsep]
    \item Inverse Problems with PINN
    \item Multi-fidelity Learning
    \item Transfer Learning in PINN
    \item AI applications in modern physics research
    \item \textbf{Final Project Start}: Team topic selection and planning
\end{itemize}

\textbf{Example Project Topics}:
\begin{enumerate}[noitemsep]
    \item Solving Schrödinger equation using PINN
    \item Electromagnetic problems with complex boundary conditions
    \item Applying AI to phase transition simulations
    \item Physics quantity prediction based on experimental data
    \item Inverse Problem solving
\end{enumerate}

\subsubsection*{Week 16: Final Project Presentation and Final Exam}
\begin{itemize}[noitemsep]
    \item Team project result presentations (20 minutes)
    \item Code review and feedback
    \item Final exam (theory + practical problems)
    \item Course summary and future learning directions
\end{itemize}

\newpage

\section{Grading Policy}

\begin{table}[h]
\centering
\begin{tabular}{@{}lll@{}}
\toprule
\textbf{Category} & \textbf{Weight} & \textbf{Details} \\ \midrule
Assignments & 30\% & 6 weekly assignments (5\% each) \\
Midterm Project & 25\% & Part I summary presentation \\
Final Project & 25\% & Team project (code + report + presentation) \\
Attendance \& Participation & 20\% & Class participation and discussion \\ \bottomrule
\end{tabular}
\end{table}

\subsection{Assignment Submission}
\begin{itemize}[noitemsep]
    \item Code submission through GitHub Repository
    \item Include execution methods and result analysis in README.md
\end{itemize}

\subsection{Final Project Requirements}
\begin{itemize}[noitemsep]
    \item \textbf{4 students per team}
    \item Collaboration through GitHub
    \item 10-15 page report (LaTeX recommended)
    \item \textbf{Team presentation}: 20-minute presentation + 10-minute Q\&A
    \item Reproducible code and data
    \item Specify each team member's contribution
\end{itemize}

\section{Lab Environment}

\subsection{Required Software}
\begin{verbatim}
# Python 3.9 or higher
conda create -n compphys python=3.10
conda activate compphys

# Required packages
pip install numpy scipy matplotlib
pip install pandas seaborn
pip install jupyter notebook
pip install torch torchvision  # PyTorch
pip install anthropic  # Claude API

# Optional packages
pip install plotly  # Interactive visualization
pip install sympy  # Symbolic computation
\end{verbatim}

\subsection{VS Code Extensions}
\begin{itemize}[noitemsep]
    \item Python
    \item GitHub Copilot (optional)
    \item LaTeX Workshop (for report writing)
\end{itemize}

\subsection{Claude AI Setup}
\begin{itemize}[noitemsep]
    \item Create Claude.ai account
    \item API Key issuance (for lab)
    \item Prompt Engineering practice
\end{itemize}

\section{Weekly Reading Materials}

\begin{description}[style=nextline,leftmargin=0pt]
\item[Week 1-2] Newman, "Computational Physics", Chapter 1-2; Python for Physics tutorial

\item[Week 3-5] MIT 6.S191 Lecture Notes 1-3; Goodfellow et al., "Deep Learning", Chapter 6

\item[Week 6-7] "Attention Is All You Need" - Vaswani et al. (2017); GPT-3 Paper (Brown et al., 2020)

\item[Week 9-12] Newman, "Computational Physics", Chapter 8 (ODEs); Landau \& Lifshitz, Classical Mechanics (reference)

\item[Week 13-14] Raissi et al., "Physics-informed neural networks" (2019); Karniadakis et al., "Physics-informed machine learning" (2021); Cuomo et al., "Scientific Machine Learning" (2022)
\end{description}

\section{Prerequisites}

\subsection*{Required}
\begin{itemize}[noitemsep]
    \item General Physics I, II
    \item Classical Mechanics
    \item Electromagnetism I
    \item Basic Python programming
\end{itemize}

\subsection*{Recommended}
\begin{itemize}[noitemsep]
    \item Quantum Mechanics I
    \item Mathematical Physics
    \item Statistical Physics
\end{itemize}

\section{Useful Links}

\subsection{Course Materials}
\begin{itemize}[noitemsep]
    \item \href{http://introtodeeplearning.com/}{MIT 6.S191 Course Website}
    \item \href{https://github.com/maziarraissi/PINNs}{Physics-Informed Neural Networks GitHub}
\end{itemize}

\subsection{Online Resources}
\begin{itemize}[noitemsep]
    \item Physics-Informed DeepONet
    \item \href{https://deepxde.readthedocs.io/}{DeepXDE Library}
\end{itemize}

\subsection{Paper Repositories}
\begin{itemize}[noitemsep]
    \item \href{https://arxiv.org/list/physics.comp-ph/recent}{ArXiv: Physics + AI}
    \item \href{https://arxiv.org/list/cs.LG/recent}{ArXiv: Machine Learning}
\end{itemize}

\section{Assignment Details}

\begin{description}[style=nextline,leftmargin=0pt]
\item[Assignment 1: Data Fitting (Week 2)] Fit experimental data with various regression models and select the optimal model

\item[Assignment 2: Optimizer Comparison (Week 4)] Compare performance of SGD, Momentum, Adam optimizers on the same problem

\item[Assignment 3: Mechanics Simulation (Week 9)] Write complex mechanics problem simulation code using LLM

\item[Assignment 4: Electromagnetic Visualization (Week 10)] Calculate electric field for given charge distribution and create visualization

\item[Assignment 5: Quantum Simulation (Week 11)] Calculate and analyze quantum states in various potentials

\item[Assignment 6: PINN Project (Week 14)] Solve selected PDE using PINN and compare with traditional methods
\end{description}

\section{Learning Outcomes}

Upon successful completion of this course, students will be able to:

\begin{enumerate}[noitemsep]
    \item Understand the mathematical principles of Neural Network operations
    \item Possess efficient coding skills using LLMs
    \item Gain practical experience in solving physics problems with AI
    \item Master differential equation solutions using PINNs
    \item Apply the latest AI technologies to physics
    \item Collaborate and manage code through GitHub
    \item Write research papers and deliver presentations
\end{enumerate}

\section{Course Philosophy}

\begin{quote}
\textit{A physicist is one who understands nature,\\
a computational physicist simulates nature,\\
and a physicist in the AI era learns and predicts nature.}
\end{quote}

This course aims not merely to teach how to use AI tools, but to learn \textbf{new problem-solving approaches by integrating physical intuition with AI technology}.

\section{Contact and Office Hours}

\begin{tabular}{ll}
\textbf{Classroom}: & TBA \\
\textbf{Lab Room}: & TBA \\
\textbf{Office Hours}: & Wednesday 14:00-16:00 (by appointment recommended) \\
\textbf{Email}: & TBA \\
\textbf{Course GitHub}: & TBA \\
\textbf{Q\&A}: & Slack channel \\
\end{tabular}

\section{Important Notes}

\begin{enumerate}[noitemsep]
    \item \textbf{Academic Integrity}: Sources must be clearly cited when using AI tools
    \item \textbf{Code Sharing}: Assignments are individual work; code copying is prohibited
    \item \textbf{Attendance}: Attendance is very important as this is a lab-oriented course
    \item \textbf{Required Materials}: Laptop required (for lab work)
\end{enumerate}

\section{Expected End-of-Semester Deliverables}

\subsection{Individual}
\begin{itemize}[noitemsep]
    \item Portfolio of 6 assignments
    \item Personal GitHub Repository
    \item Collection of physics simulation codes
\end{itemize}

\subsection{Team}
\begin{itemize}[noitemsep]
    \item Final project report
    \item Presentation materials
    \item Open-source code (GitHub)
\end{itemize}

\section{Additional Learning Resources (Optional)}

\subsection{Online Courses}
\begin{itemize}[noitemsep]
    \item Andrew Ng's Machine Learning Course (Coursera)
    \item Fast.ai Practical Deep Learning
    \item Stanford CS231n (Convolutional Neural Networks)
\end{itemize}

\subsection{YouTube Channels}
\begin{itemize}[noitemsep]
    \item 3Blue1Brown (Neural Networks Series)
    \item Two Minute Papers
    \item Arxiv Insights
\end{itemize}

\subsection{Books}
\begin{itemize}[noitemsep]
    \item \textit{Neural Networks and Deep Learning} by Michael Nielsen (free)
    \item \textit{Dive into Deep Learning} by Zhang et al. (free)
    \item \textit{Computational Physics} by Giordano \& Nakanishi
\end{itemize}

\section{Course Plan Changes}

\begin{itemize}[noitemsep]
    \item This syllabus may be adjusted according to learning progress
    \item Important changes will be announced at least one week in advance
    \item Lab content may be adjusted based on student feedback
\end{itemize}

\vspace{1cm}

\section*{Closing Remarks}

Through this course, you will master \textbf{essential AI tools for 21st-century physicists}. Beyond simply executing code, we expect you to \textbf{combine physical intuition with AI technology to achieve creative problem-solving}.

Let's make this an exciting semester together!

\vspace{1cm}

\noindent\textbf{Last Updated}: January 15, 2026\\
\textbf{Version}: 1.0\\
\textbf{Instructor}: TBA

\end{document}
